\documentclass[pmlr]{jmlr}% new name PMLR (Proceedings of Machine Learning Research)

 % The following packages will be automatically loaded:
 % amsmath, amssymb, natbib, graphicx, url, algorithm2e

 %\usepackage{rotating}% for sideways figures and tables
\usepackage{longtable}% for long tables

 % The booktabs package is used by this sample document
 % (it provides \toprule, \midrule and \bottomrule).
 % Remove the next line if you don't require it.
\usepackage{booktabs}
 % The siunitx package is used by this sample document
 % to align numbers in a column by their decimal point.
 % Remove the next line if you don't require it.
\usepackage[load-configurations=version-1]{siunitx} % newer version
 %\usepackage{siunitx}

 % The following command is just for this sample document:
\newcommand{\cs}[1]{\texttt{\char`\\#1}}

 % Define an unnumbered theorem just for this sample document:
\theorembodyfont{\upshape}
\theoremheaderfont{\scshape}
\theorempostheader{:}
\theoremsep{\newline}
\newtheorem*{note}{Note}

 % change the arguments, as appropriate, in the following:
\jmlrvolume{1}
\jmlryear{2010}
\jmlrworkshop{Workshop Title}

\title[Short Title]{Full Title of Article\titlebreak This Title Has
A Line Break\titletag{\thanks{sample footnote}}}

 % Use \Name{Author Name} to specify the name.

 % Spaces are used to separate forenames from the surname so that
 % the surnames can be picked up for the page header and copyright footer.
 
 % If the surname contains spaces, enclose the surname
 % in braces, e.g. \Name{John {Smith Jones}} similarly
 % if the name has a "von" part, e.g \Name{Jane {de Winter}}.
 % If the first letter in the forenames is a diacritic
 % enclose the diacritic in braces, e.g. \Name{{\'E}louise Smith}

 % *** Make sure there's no spurious space before \nametag ***

 % Two authors with the same address
  \author{\Name{Author Name1\nametag{\thanks{with a note}}} \Email{abc@sample.com}\and
   \Name{Author Name2} \Email{xyz@sample.com}\\
   \addr Address}

 % Three or more authors with the same address:
 \author{\Name{Laura van der Hoef} \Email{l.g.s.vanderhoef@students.uu.nl}\\
  \Name{Author Name2} \Email{an2@sample.com}\\
  \Name{Author Name3} \Email{an3@sample.com}\\
  \addr Address}



\begin{document}

\maketitle

\begin{abstract}
This is the abstract for this article.
\end{abstract}
\begin{keywords}
List of keywords
\end{keywords}

\section{Introduction}
\label{sec:intro}

\section{Data}
\label{sec:data}

\section{Data Preprocessing}
\label{sec:preprocessing}


\section{Experimental Setup}
\label{sec:setup}

All models in this experiment used a bag-of-words representation. 
This approach provides a method for representing text data by treating word counts as primary features, although it inherently lacks the ability to capture word order.
Each model was tuned using a grid search with 5-fold cross-validation to identify the optimal hyperparameters (\tableref{tab:hyperparameters}).
5-fold cross-validation was chosen to balance the need for robust evaluation against the limitations of the dataset size and available computational time.

\begin{table}[htbp]
\floatconts
  {tab:hyperparameters}%
  {\caption{Hyperparameter grids searched with 5-fold cross-validation and the selected values. Parameter names follow scikit-learn; min\_df not listed means the default of 1 was used.}}%
  {\begin{tabular}{lllc}
  \toprule
  \bfseries Model & \bfseries Hyperparameter & \bfseries Grid & \bfseries Selected\\
  \midrule
  Naive Bayes & min\_df & $\{1, 2, 3, 5, 10\}$ & 5\\
   & $k$ & $\{50, 100, 200, 300, 500, \text{all}\}$ & all\\
  \midrule
  Logistic Regression & min\_df & $\{1, 2, 3, 5, 10\}$ & 3\\
   & $C$ & 15 log-spaced in $[0.01, 100]$ & 1.931\\
  \midrule
  Decision Tree & ccp\_alpha & all $\alpha$ on the pruning path & 0.0044\\
  \midrule
  Random Forest & min\_df & $\{1, 2, 5\}$ & 2\\
   & max\_features & $\{\sqrt{p}, \log_2 p, 0.05p, 0.1p\}$ & $\log_2 p$\\
   & n\_estimators & $\{100, 300, 500\}$ & 500\\
  \midrule
  Gradient Boosting & learning\_rate & $\{0.05, 0.1, 0.2\}$ & 0.2\\
   & n\_estimators & $\{100, 300\}$ & 300\\
   & max\_depth & $\{1, 3, 5\}$ & 3\\
  \bottomrule
  \end{tabular}}
\end{table}

The first model chosen was Multinomial Naive Bayes.
As a standard model for text classification, it counts word frequencies per class.
By beginning with this baseline, we can evaluate whether increased model complexity yields improvements in accuracy.
We used the chi-squared test to select the top $k$ features, which assists in removing features that are less informative for prediction.
Additionally, the minimum document frequency ($\text{min\_df}$) was tuned via grid search to remove sparse terms.
However, the selected $k$ was all features, indicating that feature selection did not improve performance.

The second model selected was Logistic Regression.
This is another common and computationally efficient baseline for text classification, as it calculates a weight for each term within the bag-of-words representation, making the model highly interpretable.
To prevent overfitting given the limited sample size, L1 regularization was applied to penalize large weights.
The regularization strength $C$ (where $C = 1/\lambda$) was tuned over a range of 0.01 to 100.
While Multinomial Naive Bayes uses top $k$ feature selection as a tunable step, L1 regularization removes uninformative terms by setting their weights to zero, which is not possible for Naive Bayes as it has no learned weights to penalize.
As with the previous model, $\text{min\_df}$ was also employed to eliminate sparse terms.

The third model deployed was the Decision Tree.
This model was chosen for its interpretability, allowing for a clear analysis of the factors driving its predictions.
Furthermore, since a Random Forest model is also implemented, a single Decision tree serves as a necessary baseline to determine if an ensemble approach improves performance, as documented in the literature \citep{Breiman2001}.
The Decision Tree Classifier was implemented using the scikit-learn library with the default Gini setting.
However, since the default $\text{ccp\_alpha}$ is 0, we performed a grid search to optimize this hyperparameter.
Without tuning $\text{ccp\_alpha}$ the tree can become overly complex and lead to overfitting.
The final model for testing was refit on the entire training set once the optimal $\alpha$ was found. 



% Ensemble models

In addition to accuracy, precision, recall and F1-score, we used McNemar's test \citep{Dietterich1998} to test whether difference in accuracy between the models is significant.
McNemar's test is computationally efficient and requires only a single train/test split, suiting our setup.
We compared Naive Bayes and Logistic Regression because they represent the generative and discriminative models, respectively \citep{Ng2001}.
The results of this test will indicate whether for this text classification task it is more effective to model the estimated generative probability of the words or to model the class boundary more directly.
% Motivate McNemar comparison of other models here (Q2)

We used the McNemar formula as follows:
\begin{equation}\label{eq:mcnemar}
\chi^2 = \frac{\left(|n_{01} - n_{10}| - 1\right)^2}{n_{01} + n_{10}}
\end{equation}
This formula uses a $2 \times 2$ contingency table that counts the instances where the two models agree or disagree, where $n_{01}$ is the number of instances misclassified by the first model but not the second, and $n_{10}$ vice versa.
The null hypothesis states that both models perform with equal accuracy, while the alternative hypothesis states that the models differ in accuracy, tested at a significance level of $\alpha = 0.01$.


\section{Results}
\label{sec:results}

\subsection{Performance}
\label{sec:performance}

\begin{table}[htbp]
\floatconts
  {tab:performance}%
  {\caption{Cross-validation accuracy on the training set and test set performance of the final models.}}%
  {\begin{tabular}{lSSSSS}
  \toprule
  \bfseries Model & {\bfseries CV accuracy} & {\bfseries Accuracy} & {\bfseries Precision} & {\bfseries Recall} & {\bfseries F1}\\
  \midrule
  Naive Bayes & 0.9005 & 0.8980 & 0.8566 & 0.9560 & 0.9036\\
  Logistic Regression & 0.9071 & 0.9020 & 0.9170 & 0.8840 & 0.9002\\
  Decision Tree & 0.7889 & 0.7800 & 0.7652 & 0.8080 & 0.7860\\
  Random Forest & 0.9172 & 0.8960 & 0.9195 & 0.8680 & 0.8930\\
  Gradient Boosting & 0.9078 & 0.9000 & 0.9098 & 0.8880 & 0.8988\\
  \bottomrule
  \end{tabular}}
\end{table}

\tableref{tab:performance} shows the performance metrics for the models evaluated. 
Multinomial Naive Bayes and Logistic Regression achieved similar accuracy (0.898 and 0.902) and F1-scores (0.904 and 0.900).
However, they show differences in precision and recall.
Naive Bayes has a higher recall (0.956) but lower precision (0.857), so it misses fewer AI-generated reviews but also incorrectly labels more human reviews as AI.
Logistic Regression shows the opposite pattern with a higher precision (0.917) and lower recall (0.884).
This means it incorrectly flags fewer human reviews as AI-generated.
In comparison to the other models, the Decision Tree performs worst on all metrics.
It has an accuracy of 0.780 and an F1-score of 0.786.
It shows a similar error pattern to Naive Bayes where its recall (0.808) is higher than its precision (0.765). However, the difference is smaller in comparison.

\subsection{Statistical Tests}
\label{sec:statistical-tests}

The McNemar's test results for Multinomial Naive Bayes and Logistic Regression are detailed in \tableref{tab:mcnemar-lr-nb}.
As shown, both models correctly classify the vast majority of reviews (421 instances).
Logistic regression misclassifies 28 reviews that Naive Bayes correctly identifies, while the reverse occurs in 30 instances.
Given the proximity of these numbers, McNemar's test gives $p = 0.896 > 0.01$, so we failed to reject the null hypothesis.
This means Logistic Regression and Naive Bayes do not differ significantly in accuracy.

\begin{table}[htbp]
\floatconts
  {tab:mcnemar-lr-nb}%
  {\caption{McNemar contingency table for Logistic Regression vs.\ Naive Bayes on the test set ($\chi^2 = 0.017$, $p = 0.896$).}}%
  {\begin{tabular}{lcc}
  \toprule
  & \bfseries NB correct & \bfseries NB incorrect\\
  \midrule
  \bfseries LR correct & 421 & 30\\
  \bfseries LR incorrect & 28 & 21\\
  \bottomrule
  \end{tabular}}
\end{table}

\subsection{Feature Importance}
\label{sec:feature-importance}

\begin{table}[htbp]
\floatconts
  {tab:top-features-linear}%
  {\caption{Top 5 features of the linear models. For Naive Bayes the score is the log-probability ratio $\log P(w \mid \text{AI}) - \log P(w \mid \text{human})$; for Logistic Regression it is the L1 coefficient. Positive scores point towards AI-generated, negative towards human-written reviews.}}%
  {\begin{tabular}{lSlSlSlS}
  \toprule
  \multicolumn{4}{c}{\bfseries Naive Bayes} & \multicolumn{4}{c}{\bfseries Logistic Regression}\\
  \cmidrule(lr){1-4}\cmidrule(lr){5-8}
  \multicolumn{2}{c}{AI} & \multicolumn{2}{c}{Human} & \multicolumn{2}{c}{AI} & \multicolumn{2}{c}{Human}\\
  \midrule
  delight & 3.4233 & told & -3.8778 & delight & 3.5810 & enter & -3.7465\\
  notch & 3.2563 & said & -3.6414 & lack & 2.8383 & etc & -2.7759\\
  desir & 3.2255 & never & -3.6075 & offer & 2.8076 & big & -2.6572\\
  inconsist & 3.1938 & com & -3.4183 & stylish & 2.5321 & fan & -2.5060\\
  lastli & 3.1610 & second & -3.3313 & notch & 2.5021 & us & -2.1774\\
  \bottomrule
  \end{tabular}}
\end{table}

The top five most important features for Multinomial Naive Bayes (\tableref{tab:top-features-linear}) are: \emph{delight}, \emph{notch}, \emph{desir}, \emph{inconsist}, and \emph{lastli}.
All five features show similar feature scores, suggesting the model assigns them nearly equal importance.
This means each word was approximately $e^{3.16} \approx 24$ to $e^{3.42} \approx 31$ times more likely to appear in AI-generated reviews.
These terms appears with relatively high frequency in the dataset (around 30 to 40 occurrences each), implying that their high occurrence in AI reviews was a primary driver for the Naive Bayes classification.
The words associated with AI-generation show formal characteristic, such as ``delight'' and ``notch,'', suggesting that GPT-4's output is more formal compared to human writing.
``Inconsis'' was also a top feature for AI, likely derived from ``inconsistent'' or ``inconsistently''. 
``Delight'' likely stems from ``delightful'' or ``delighted'', while ``notch'' likely from ``top-notch''.
Furthermore, ``inconsist'' and ``notch'' are primarily adjectives.
This suggests that GPT-4 relies on descriptive adjectives.
``Lastly'' is a distinct outlier in the list of words, since it is a transition word.
The use of it suggests that GPT-4 uses words that are structural connectors more frequently than humans.

The top five features for the Logistic Regression model were: \emph{delight}, \emph{lack}, \emph{offer}, \emph{stylish}, and \emph{notch}.
The two words \emph{delight} and \emph{notch} overlap between the two linear models, with \emph{delight} ranking highest for both.
The variance in word overlap can be attributed to Naive Bayes modeling the generative probability of individual words, whereas Logistic Regression models the class boundary.
\emph{Stylish} is an adjective, further suggesting that GPT-4 leans towards descriptive words.
Additionally, \emph{lack} was identified as a sign of AI-generation, suggesting that AI may focus more on missing services, whereas human reviews may do so less frequently.
While two AI-related words overlap between the linear models, none of the human words match.
This may be because AI produces more formulaic reviews, whereas human writing is more varied.

The features associated with AI generation appear to be more formal and polished in nature.
Furthermore, AI tends to describe the appearance of locations in polished terms, whereas humans focus more on actions (e.g. ``told,'' ``said,'' ``enter'').
This may suggest a difference in that AI describes experiences based on impressions of what is likely to occur, while humans describe their lived experiences.
One outlier ``com'' appears among the human features.
This is likely a fragment of ``booking.com'' as a result of data preprocessing. 
This resulted in Naive Bayes picking up on a platform-specific artifact of the human reviews, whereas AI does not mention the booking platforms.
A limitation of the linear models is the bag-of-words representation, which obscures the analysis of word order.



The top five features for the Decision Tree (\tableref{tab:top-features-trees}), ranked by Gini importance, are: \emph{servic}, \emph{staff}, \emph{room}, \emph{lack}, and \emph{good}.
These features show no overlap with the linear models.
However, unlike the previous models, the AI-associated words are mainly about the hotel itself, such as \emph{servic} (service), \emph{staff} and \emph{room}, likely because the tree favours common words that split many reviews.
The only exception is \emph{good}, which is associated with human reviews.
This word also has a significantly lower Gini importance than the other top four features.
Although descriptive words are mostly associated with AI in the other models, \emph{good} points towards human reviews. This could be because AI uses more polished language (e.g., \emph{delight}, \emph{top-notch}, \emph{stylish}) to appear convincing, whereas humans may use simpler, more neutral adjectives like \emph{good} to express satisfaction.


\begin{table}[htbp]
\floatconts
  {tab:top-features-trees}%
  {\caption{Top 5 features of the tree-based models, ranked by Gini importance. Direction is the class with the larger share of training reviews containing the word.}}%
  {\begin{tabular}{lSclSclSc}
  \toprule
  \multicolumn{3}{c}{\bfseries Decision Tree} & \multicolumn{3}{c}{\bfseries Random Forest} & \multicolumn{3}{c}{\bfseries Gradient Boosting}\\
  \midrule
  servic & 0.2629 & AI & room & 0.0209 & AI & servic & 0.1047 & AI\\
  staff & 0.1899 & AI & servic & 0.0201 & AI & room & 0.0921 & AI\\
  room & 0.0998 & AI & staff & 0.0131 & AI & staff & 0.0754 & AI\\
  lack & 0.0765 & AI & hotel & 0.0128 & AI & lack & 0.0352 & AI\\
  good & 0.0561 & Human & lack & 0.0109 & AI & central & 0.0314 & AI\\
  \bottomrule
  \end{tabular}}
\end{table}



\section{Conclusions}
\label{sec:conclusions}


\bibliography{references}

\section*{AI Statement}

\paragraph{Laura van der Hoef}

I used Claude to review my code, suggest readability improvements and help solve bugs. For example, it suggested how to format the log statements so the output was easier to read and interpret.
I also used Claude to help me with styling in Latex, such as setting up the tables in the results section and formatting the text correctly.
Claude was used for finding spelling \& grammar mistakes and providing suggestions to fix them.

\appendix

\section{First Appendix}\label{apd:first}

This is the first appendix.

\section{Second Appendix}\label{apd:second}

This is the second appendix.

\end{document}
